\section{Introducci\'on}

Imagina que trabajas para Shosholanda y, de buenas a primeras, te enteras de que dos de sus
agentes, Miri y Chipotles, andan soltando contrase\~nas de la empresa a gente de fuera.
El CEO, Shiro, nos contrat\'o para investigar el caso, pero con una condici\'on: nada de
preguntarles directamente, porque si se dan cuenta, en un parpadeo borrar\'ian todo rastro.
As\'i que la \'unica herramienta que tenemos es lo que ellos mismos dejaron a la vista.

A eso se le llama OSINT (\textit{Open Source INTelligence}): buscar, juntar y analizar informaci\'on
que es p\'ublica. No siempre est\'a en internet, tambi\'en puede salir de noticias, registros
p\'ublicos o hasta de lo que cuentan los conocidos. Lo curioso es que casi nadie se da cuenta
de cu\'anto deja en l\'inea: lo que nos gusta, d\'onde andamos, con qui\'en andamos y hasta
c\'omo pensamos. Si esa informaci\'on cae en manos con malas intenciones, alguien podr\'ia armar
un ataque hecho a la medida; y en el mejor de los casos, solo nos conquista el crush.

Esta pr\'actica tiene tres metas: (1) tomar conciencia de lo que publicamos y de que el alcance
es, potencialmente, todo el mundo; (2) entender c\'omo los metadatos de un archivo pueden decir
mucho m\'as de lo que se ve a simple vista; y (3) conectar eventos para reconstruir historias.
Todo se hizo con fuentes abiertas, con fines did\'acticos y sin intenci\'on de doxxear a nadie.
